\subsection{Local multi-root cancellation and the remaining common-scale problem}
\label{ssec:multiroot-filter}

\paragraph{Status (10 September 2026).}
The identities below are exact finite statements under their explicit hypotheses.
The reported evaluations use point arithmetic, not outward-rounded numerical
certificates. They improve the a posteriori transfer from a finite eigenstate
to a root displacement; they do not supply the missing a priori law in
$c=\lambda^2,N$ and precision. In particular, no fitted conversion constant or
reference-root error enters the new predictor.

Retain the normalization and observables of
\eqref{eq:divided-observables}. For a stencil $\mathcal J$ of $p+1$ distinct
simple positive nonlattice finite roots $\nu_i$, set
\[
 t_i=\nu_i^2,\qquad a_i^*=\frac{2t_i}{\kappa_i},\qquad
 y_i=\frac{J_i-f_i^2-\varepsilon_N h_iS_i}
                {\varepsilon_N\kappa_i}.
\]
Assume $\varepsilon_N f_i\kappa_i\ne0$ and correctly match each root to a
distinct simple real ordinate $\gamma_i$. Put $d_i=\gamma_i-\nu_i$,
$B_i=F(\gamma_i)/d_i$, continued by $f_i$, and define
\[
 L_i=\frac{2\nu_i d_i B_i^2}{2\nu_i+d_i},\quad
 Q_i=B_i^2\left(1+\frac{d_i^2}{(2\nu_i+d_i)^2}\right),\quad
 C_i=\frac{Q_i-f_i^2-h_iL_i}{\varepsilon_N\kappa_i}.
\]
Exclude every stencil pair from a \emph{common} remote multiset, write
$x_m=z_m^2$, $E_m=2F(z_m)^2$ for positive-real-part representatives, including
multiplicity, and set $M(t)=\sum_{\rm common}(E_m/\varepsilon_N)/(x_m-t)$.
Assume $\gamma_j^2\ne t_i$ for $i\ne j$ and $x_m\ne t_i$ in the common
remote set. Off-line summands retain complex squares and conjugate pairing.
The finite Weil identities give exactly
\begin{equation}\label{eq:multiroot-robin}
 y_i=M(t_i)+a_i^*M'(t_i)+C_i+
 \sum_{j\in\mathcal J\setminus\{i\}}\frac{E_j}{\varepsilon_N}
 \left(\frac1{\gamma_j^2-t_i}+\frac{a_i^*}{(\gamma_j^2-t_i)^2}\right).
\end{equation}
The local pair contributes $L_i$ to $\varepsilon_NS_i$ and $Q_i$ to $J_i$.
The respective remote-pair contributions are
$\nu_i E_m/(x_m-t_i)$ and $E_m(x_m+t_i)/(x_m-t_i)^2$.
Subtraction uses $1-\nu_i h_i=\kappa_i$.

\begin{proposition}[Polynomial-reproducing remote cancellation]
\label{prop:multiroot-cancellation}
Fix a target $0\in\mathcal J$. Suppose weights $w_i$ solve
\begin{equation}\label{eq:multiroot-weights}
 \sum_{i\in\mathcal J}w_i\bigl[P(t_i)+a_i^*P'(t_i)\bigr]=P(t_0)
 \quad(\deg P\le p).
\end{equation}
Then the operator-derived predictor
\begin{equation}\label{eq:multiroot-predictor}
 \widehat\delta_p=\frac{\varepsilon_N}{f_0^2}
              \left(\nu_0\sum_i w_i y_i-S_0\right)
\end{equation}
has remote error kernel
\begin{equation}\label{eq:multiroot-kernel}
 K(x)=\sum_i w_i\left(\frac1{x-t_i}+\frac{a_i^*}{(x-t_i)^2}\right)
                  -\frac1{x-t_0}=O(x^{-p-2}).
\end{equation}
For $p=0$ this is the previous divided-form correction.
\end{proposition}
\begin{proof}
Expand each rational term at infinity. Its coefficient of $x^{-q-1}$ is
$t_i^q+q a_i^*t_i^{q-1}$. Equation~\eqref{eq:multiroot-weights} makes this
coefficient equal to $t_0^q$ for $q=0,\ldots,p$. This proves the stated
cancellation, without assuming that the remote transform is a polynomial.
\end{proof}

For completeness define
\[
 X=\sum_{j\in\mathcal J}\frac{E_j}{\varepsilon_N}
 \left\{\nu_0\sum_{i\ne j}w_i
 \left[\frac1{\gamma_j^2-t_i}+\frac{a_i^*}{(\gamma_j^2-t_i)^2}\right]
       -\frac{\mathbf1_{j\ne0}\nu_0}{\gamma_j^2-t_0}\right\}.
\]
Direct substitution in~\eqref{eq:multiroot-robin} proves the exact remainder
\begin{equation}\label{eq:multiroot-remainder}
 \widehat\delta_p-(\nu_0-\gamma_0)=
 \frac{\varepsilon_N}{f_0^2}
 \left\{\nu_0\sum_{\rm common}\frac{E_m}{\varepsilon_N}K(x_m)
                +X+\nu_0\sum_iw_iC_i\right\}
                  +d_0-\frac{L_0}{f_0^2}.
\end{equation}
In particular, the cross-stencil and nonlinear local terms cannot be omitted
from a proof merely because they are small in a numerical example.

\paragraph{An explicit finite-height bound.}
Let $\Delta_i=t_i-t_0$, $\max_i\nu_i<H<T$, and $T>\alpha_N$. Require every matched stencil ordinate to lie below $H$. Define
\begin{align*}
 B_K(H)=\sum_i\frac{|w_i|}{(H^2-t_0)^{p+1}}
 \biggl\{&\frac{|\Delta_i|^{p+1}}{H^2-t_i}\\
 &+|a_i^*|\left[\frac{(p+1)|\Delta_i|^p}{H^2-t_i}
             +\frac{|\Delta_i|^{p+1}}{(H^2-t_i)^2}\right]\biggr\}.
\end{align*}
Expanding $1/(x-t_i)$ to order $p$ in $\Delta_i$, with exact remainder
$\Delta_i^{p+1}/[(x-t_0)^{p+1}(x-t_i)]$, and differentiating that remainder
proves $|K(z^2)|\le B_K(H)$ when $|\Re z|\ge H$.
Here $|z^2-t_i|\ge H^2-t_i$, including off-line zeros.
Platt--Trudgian's finite verification permits $T=3\cdot10^{12}$
\cite{AttackPlattTrudgian}. For a complete pair-energy prefix $E_H$ and an
absolute complex-tail bound $V_T$, the uncomputed portion of
$\sum_{\rm common}E_mK(x_m)$ is bounded by
\begin{equation}\label{eq:multiroot-tail}
 (\varepsilon_N-E_H+V_T)B_K(H)+V_TB_K(T).
\end{equation}
Positivity is used only between $H$ and $T$; the rest uses absolute variation.
A valid conservative choice from the finite rational expansion is
\begin{gather*}
 V_T\le\frac{16}{L}\cosh^2(L/4)B_M(T)^2\frac{\log T+1}{T},
 \qquad \mu_{2m}=\sum_j a_j\alpha_j^{2m},\\
 B_M(T)=\sum_{m<M}\frac{|\mu_{2m}|}{T^{2m}}+
 \frac{\|a\|_1(\alpha_N/T)^{2M}}{1-(\alpha_N/T)^2}.
\end{gather*}
Here $M\ge1$ is an integer and $\alpha_N=2\pi N/L$.
Corollary~1.2 of Hasanalizade--Shen--Wong \cite{AttackZeroCount}
implies $\mathcal N(t)\le t\log t$ for $t\ge T$. Integration by parts gives
$\sum_{\Im\rho>T}(\Im\rho)^{-2}\le2(\log T+1)/T$; the finite even
rational expansion and $|\sin(Lz/2)|\le\cosh(L/4)$ in the critical strip
then give the displayed variation bound. No assumption of global RH is used.

If $|d_i|\le r_i<2\nu_i$ and $|F'''|\le M_{3,i}$ on the corresponding real
interval, put $b=M_{3,i}/(6|f_i|)$,
\begin{gather*}
 C=2b+(|h_i|/2+br_i)^2,\qquad
 G=\frac{|h_i|}{2\nu_i-r_i}+\frac1{(2\nu_i-r_i)^2},\\
 H_i=C+G+(|h_i|+Cr_i)(|h_i|+Gr_i),\\
 R_i=(1+|h_i|r_i/2+br_i^2)^2.
\end{gather*}
Taylor's theorem, with the first-order terms canceled, gives
\[
 |C_i|\le\frac{f_i^2r_i^2H_i}{|\varepsilon_N\kappa_i|},\qquad
 |d_i-L_i/f_i^2|\le r_i^2
       \left(|h_i|+Cr_i+\frac{R_i}{2\nu_i-r_i}\right).
\]
Together these formulas give a fully specified conditional finite enclosure.
Matching, certified coefficients and directed rounding are separate obligations.
If the remainder is enclosed by $[A,B]$, the signed relative displacement
error is obtained from $q/(\widehat\delta_p-q)$ at the endpoints only when
$\widehat\delta_p\notin[A,B]$. An interval containing that pole cannot be
reported as a finite relative-error enclosure; it must remain unassessed.

\paragraph{Conditioning and falsification.}
Simplicity alone does not imply $\kappa_i\ne0$: the even polynomial
$F(z)=z^2-4$ at $\nu=2$ is an algebraic counterexample. Nor do nonzero kappas
imply an invertible stencil: $F=(z^2-1)(z^2-4)$ gives
$t_i+a_i^*=5/2$ at both roots, making the $p=1$ system singular. These are
algebraic fixtures, not asserted CCM ground-state examples. A CCM-specific
uniform conditioning theorem remains open.

The frozen primary choice was $p=2$, using the target and the two nearest
available computed roots by $|\nu_i-\nu_0|$, ties to lower $\nu_i$.
It passed the $1\%$ displacement gate on 110 retained configuration/root pairs.
Point substitution into the finite-height remainder covered all 90 pairs with
complete energy/Taylor inputs; 20 others lacked those audit inputs. In the
original sparse $c=1000,N=890$ pool, the signed magnitude errors at $k=600,610$
were $-0.4876697\%$ and $+0.5583456\%$, with respective point ranges
$[-0.4876807,-0.4876566]\%$ and $[0.5583284,0.5583617]\%$.
A separately frozen acquisition of 11 additional indices gave a denser 41-root
pool. All 41 primary predictions passed; the errors at $k=595,600,610$ became
$+0.0559910\%$, $-0.00314985\%$, $+0.00447358\%$. The previous one-root
correction fails at $k=595$ as well, by $-1.8918383\%$.
The new indices use reference-seeded, full-precision Newton validation on the
same retained eigenstate, not independent discovery or 41 independent assemblies.
There are 121 distinct configuration/root pairs across the two stages,
because 30 of the 41 dense-pool cases repeat earlier indices.
The fresh dense cases are not included in the earlier 90-case remainder audit.
The largest weight $1$-norm in the dense primary rule is $4.408$; at degree
$4$ it exceeds $1156$. Higher order therefore requires a conditioning audit,
even when its point score improves. Rounding the initial predictor inputs
$y_i,\kappa_i$ to 20 significant digits changes its value by at most
$1.11\cdot10^{-20}$ relatively in 110 tests; this checks postprocessing
sensitivity, not eigenstate accuracy.
Global Pad\'e and nonlinear rational alternatives failed near the edge and are
retained as failures in the research return packet. The previously frozen
shifted-erf eigenvalue model also fails the denser $c=1200,N=60$--$150$ window
($1.166\%$ RMS and $1.424\%$ maximum depth error) and gives the wrong cutoff-derivative
sign at $c=100,N=300$. These are dated evidence against that model, not a
contradiction of an unrelated coarse value envelope or a new fit.

\paragraph{The exact polynomial window and the common factor.}
Let $\beta=2\pi/L$, $x_j=\beta^2j^2$, $w_0=1$, $w_j=\sqrt2$ for $j>0$,
$Q_N(x)=\prod_{j=0}^N(x-x_j)$ and
$P_u(x)=Q_N(x)\sum_{j=0}^N w_ju_j/(x-x_j)$. Then
\begin{equation}\label{eq:attack-gamma-bridge}
 F_u(z)=\frac{(-1)^N2\pi}{\sqrt L\,\beta^{2N+1}}
       \frac{P_u(z^2)}{\Gamma(N+1+z/\beta)\Gamma(N+1-z/\beta)}.
\end{equation}
This follows from the finite Gamma product and sine reflection identities,
including removable singularities by reciprocal Gamma. Thus at most $N$
positive nonlattice roots are movable polynomial roots; reciprocal Gamma
supplies forced lattice zeros beyond the Fourier band. A consecutive prefix
of $k$ distinct movable positive roots therefore requires $k\le N$. An
isolated root matched to a higher ordinal does not imply that every preceding
root has been recovered. This is a counting restriction, not a matching theorem.

For integers $N\ge1$, real $|\tau|<1$, and $\Phi_N(s)=\Gamma(N+1)^2/[\Gamma(N+1+s)\Gamma(N+1-s)]$,
positive-real Stirling remainders \cite{AttackDLMF} prove
\begin{gather*}
 \left|\log\Phi_N(N\tau)+NV(\tau)+\tfrac12\log(1-\tau^2)\right|
 \le\frac1{6N(1-\tau^2)},\\
 V(\tau)=(1+\tau)\log(1+\tau)+(1-\tau)\log(1-\tau).
\end{gather*}
In particular $V=\tau^2+\tau^4/6+\cdots$; a Gaussian replacement is only
central-window information, not a band-wide approximation.
The exact extremal quotient uses
\begin{gather*}
 W_N(z)=\frac{2L}{\Gamma(N+1+z/\beta)^2\Gamma(N+1-z/\beta)^2},\\
 v_0=(N!)^{-4},\qquad v_j=\frac{2}{[(N-j)!(N+j)!]^2}\quad(j>0).
\end{gather*}
These are the numerator weights on positive-real-part zero representatives
and the denominator weights on $x_j$, respectively. Then
the precise finite problem is
\[
 \varepsilon_N=\min_{0\ne P\in\mathcal P_N}
 \frac{\sum_{\Re z>0}W_N(z)P(z^2)^2}{\sum_{j=0}^Nv_jP(x_j)^2}.
\]
Here the numerator includes multiplicities and complex squares, and defines
the positive finite polynomial inner product inherited from the exact even
matrix. For its reproducing kernel $K_N$, the elementary Gram-inverse
and Christoffel principles \cite{AttackSimon} give
\[
 T_1=\sum_jv_jK_N(x_j,x_j),\qquad
 T_2=\sum_{i,j}v_iv_jK_N(x_i,x_j)^2.
\]
The signed-complex numerator must be retained without global RH; finite
positivity suffices for this finite-dimensional kernel identity.
A uniform relative estimate of this arithmetic weighted kernel is still missing.

\paragraph{Residual and precision obligations.}
The divided-form algebra specifies errors that an implementation must retain.
In full Fourier coordinates write $D=\operatorname{diag}(\alpha_j)$,
$[D,A]=g e^T-e g^T$ with fixed boundary vector $e=(1,\ldots,1)^T$,
$b=(D-\nu I)^{-1}a$, $q=Aa-\varepsilon_Na$,
$r=-e^Tb=\sum_j a_j/(\nu-\alpha_j)$, $G=g^Tb$, and
$H=g^T(D-\nu I)^{-1}b$. Then $r'=-e^T(D-\nu I)^{-1}b$
and the displacement identity gives
\[
 b^TAb=\varepsilon_N\|b\|^2+G r'-rH
                  +b^T(D-\nu I)^{-1}q.
\]
The root-residual term $-rH$ is retained in the new dense postprocessing;
the stored eigenvector alone does not certify the last term.

For normalized, sign-matched finite states with $\|u-\widetilde u\|\le\eta_u$,
the Fourier integral and Cauchy--Schwarz give the real-axis bound
\[
 |F_u^{(m)}-F_{\widetilde u}^{(m)}|
 \le \frac{\eta_u L^{m+1/2}}{2^m\sqrt{2m+1}}.
\]
If $F_u(\nu)=0$ and $|F_u'|\ge m_0>0$ on $[\nu-r,\nu+r]$,
$\eta_u\sqrt L<m_0r$, and $\eta_u L^{3/2}/\sqrt{12}<m_0$, the perturbed
transform has a unique zero in that interval, displaced by at most
$\eta_u\sqrt L/m_0$. This is a conditional stability bound. Certified state
error, root matching, derivative lower bounds, and stencil conditioning are
still needed for an end-to-end precision guarantee.

\paragraph{The arithmetic innovation profile.}
The equivalent computational route uses the even displacement nodes $j^2$ and
scalar Schur jets: for a pivot $\sigma_N(z)$ of $A_N-zI$,
\[
 I_N=-\sigma_N'(0)/\sigma_N(0),\qquad
 T_p=-p[z^p]\sum_{j=0}^N\log\sigma_j(z).
\]
Fixed-order rank-two generator elimination costs $O(N^2)$ scalar operations.
Its exact recurrence was checked on 1,032 rational identities, and an independent
analytic assembly reproduced all 366 retained $c=500$ prefix exports of
$\sigma,M,T_1,T_2$ to their full 80 significant digits. This does not constitute
a derived innovation asymptotic or a stability certificate. At fixed $c$, if
the continuum bottom is positive, logarithmic diagonal growth instead implies
$T_p(N)\sim N/(\log N)^p$ for each fixed integer $p\ge1$; eventual fixed-order ground-state
dominance is impossible. The present ramp and the unlimited late-$N$ regime
therefore require different uniformity statements.

The boundary-column-only closure was also tested without fitted constants on
every $N=320,\ldots,1600$ at $c=5,9$. Only one of 1281 rows at each cutoff
passed a $1\%$ gate. The precise profile remains an open arithmetic problem;
it cannot be replaced by that failed local boundary approximation.
At a prime-power event $c=q=p^r$, the one-sided cutoff derivatives obey
\[
 \operatorname{jump}\frac{d[-\log_{10}\varepsilon_N]}{dL}
 =\frac{2\log p}{\sqrt q\log q}\,
   \frac{(w^Tu)^2}{\varepsilon_N\log 10},
 \qquad w=(1,\sqrt2,\ldots,\sqrt2)^T.
\]
Hellmann--Feynman gives this identity from the rank-one arithmetic jump.
It constrains a uniform derivative law;
it does not by itself falsify a coarse approximation to eigenvalue values.
